% ============================================================
% CAPITOLUL 4 — ARHITECTURA SISTEMULUI E-MEDICINE
% ============================================================

\chapter{Arhitectura Sistemului E-Medicine}
\label{ch:arhitectura}

\section{Vedere de ansamblu}
\label{sec:arhitectura_overview}

\todosec{Diagramă arhitecturală globală a sistemului: cluster K8s cu YugabyteDB,
runner-e SIR (R, Julia, Python, SAS), API Go, frontend Vue.js. Descrie fluxul
de date de la ingestia datelor epidemiologice până la vizualizare.}

\section{Cerințe funcționale și non-funcționale}
\label{sec:cerinte}

\subsection{Cerințe funcționale}
\label{subsec:cerinte_functionale}

\todosec{Lista completă de cerințe funcționale (use cases):
(1) Încărcare seturi de date epidemiologice,
(2) Declanșare rulare model SIR într-un limbaj ales,
(3) Vizualizare rezultate și grafice,
(4) Comparare rezultate între limbaje,
(5) Export rezultate (CSV, PDF).}

\subsection{Cerințe non-funcționale}
\label{subsec:cerinte_nonfunctionale}

\todosec{Cerințe non-funcționale: performanță (latența API), disponibilitate
(SLA YugabyteDB pe K8s), scalabilitate (număr simultan de rulări), securitate
(autentificare JWT), extensibilitate (adăugare noi limbaje).}

\section{Arhitectura bazei de date distribuite}
\label{sec:arhitectura_bd}

\subsection{Schema relațională}
\label{subsec:schema_bd}

\todosec{Diagrama ERD completă. Tabelele principale:
\texttt{datasets} (metadate seturi de date),
\texttt{epidemiological\_records} (valori S/I/R per timestamp și regiune),
\texttt{simulation\_runs} (parametri și metadate ale rulărilor),
\texttt{simulation\_results} (rezultatele S(t), I(t), R(t) per rulare),
\texttt{benchmarks} (timpii de execuție per configurație hardware),
\texttt{languages} (R, Julia, Python, SAS).}

\subsection{Modelul datelor epidemiologice}
\label{subsec:model_date}

\todosec{Descrierea schemei de date pentru seria temporală epidemiologică:
granularitate (zilnică/săptămânală), acoperire geografică, câmpuri obligatorii
vs. opționale. Justificare alegeri de indexare.}

\subsection{Strategia de partitionare și replicare în YugabyteDB}
\label{subsec:strategie_partitionare}

\todosec{Range sharding pe (disease\_code, date) pentru interogări de serii
temporale. Replicare Raft cu factor 3. Configurarea tablet-elor și tuning
al compaction-ului RocksDB.}

\section{Arhitectura API-ului backend Go}
\label{sec:arhitectura_api}

\subsection{Structura modulară a proiectului}
\label{subsec:structura_go}

\todosec{Structura directorului Go: cmd/, internal/(api, db, runner, models),
pkg/(sirmath). Descriere pattern-uri arhitecturale: repository pattern, service
layer, dependency injection.}

\subsection{Diagrama endpointurilor REST}
\label{subsec:endpoints}

\todosec{Tabel complet al endpointurilor REST:
GET /api/datasets, POST /api/datasets,
POST /api/simulations (trigger rulare),
GET /api/simulations/{id}/results,
GET /api/benchmarks,
GET /api/health. Documentație OpenAPI.}

\subsection{Integrarea cu runner-ele de algoritmi}
\label{subsec:integrare_runnere}

\todosec{Cum apelează API-ul Go implementările SIR: subprocess cu stdin/stdout
JSON, sau binare native? Gestiunea asincronă a rulărilor (goroutines, job queue).}

\section{Arhitectura frontend Vue.js}
\label{sec:arhitectura_frontend}

\subsection{Structura componentelor}
\label{subsec:structura_vue}

\todosec{Arborele de componente Vue: layout-uri, pagini (Dashboard, Datasets,
Simulations, Benchmarks, Compare), componente reutilizabile (SIRChart,
BenchmarkTable, ParameterForm).}

\subsection{Fluxul de date și state management}
\label{subsec:state_management}

\todosec{Diagrama fluxului de date: Pinia stores (datasetsStore, simulationsStore,
benchmarkStore), apeluri Axios la API-ul Go, reactive updates în componente.}

\section{Infrastructura Kubernetes}
\label{sec:infrastructura_k8s}

\subsection{Topologia cluster-ului}
\label{subsec:topologie_k8s}

\todosec{Diagrama cluster-ului K8s: master node(s), worker nodes, noduri
YugabyteDB (StatefulSet x3). Specifică configurația hardware minimă recomandată.}

\subsection{Deployment YugabyteDB pe Kubernetes}
\label{subsec:deploy_yugabyte}

\todosec{Helm chart YugabyteDB, StatefulSet configuration, PersistentVolumeClaims,
Services (headless service pentru DNS intern, LoadBalancer pentru accesul extern).}

\subsection{Deployment API Go și frontend Vue.js}
\label{subsec:deploy_api_frontend}

\todosec{Deployment și Service pentru API Go (containerizat cu Dockerfile multi-stage),
Deployment Nginx pentru frontend Vue.js, Ingress controller (nginx-ingress sau
Traefik) pentru rutarea traficului extern.}
